\documentclass[12pt]{article}
\usepackage[T2A]{fontenc}
\usepackage[utf8]{inputenc}
\usepackage[russian]{babel}
\usepackage{latexsym,amssymb,amsthm}
\usepackage{amsmath}
\usepackage{wrapfig}
\RequirePackage[pdftex]{graphicx}
%%\usepackage{array}

\usepackage{epsf,epic,eepic}
\input ris.tex
% \input cells.sty
% \addrisoption{\risleftfalse}

\DeclareRobustCommand{\No}{\ifmmode{\nfss@text{\textnumero}}\else\textnumero\fi}

\oddsidemargin=0pt
\textwidth=159mm
\headheight=0pt
\headsep=0pt
\topmargin=0pt
\textheight=240mm

%%
%% Три точки для обозначения делимости
%%
\catcode`\@=11
\def\kratno{\mathrel{\smash{\lower.5ex\hbox{$\,\vdots\,$}}}}
% Это прямые три точки для обозначения делимости
\def\nekratno{\mathrel{\mathpalette\c@ncel\kratno}}
\def\c@ncel#1#2{\m@th\ooalign{$\hfil#1\mkern1mu/\hfil$\crcr$#1#2$}}
% Это перечеркнутые три точки для обозначения неделимости
\def\divis{\smash{\lower.1ex\hbox{\,\hbox{.}\kern-.22em\lower-.6ex\hbox{.}
\kern-.52em\lower-1.2ex\hbox{.}\,}} }
% А это наклонные...

\hyphenation{чис-ло чис-ла чис-лу чис-лом чис-ле чис-лам чис-ла-ми чис-лах
че-ты-рех-уголь-ник че-ты-рех-уголь-ни-ка пря-мо-уголь-ник тре-уголь-ни-ка
тре-уголь-ни-ков вне-впи-сан-ной вне-впи-сан-ных}

\sloppy
\pagestyle{empty}

\def\q#1.{{\bf #1. }}
\def\dotline{\smallskip\hbox to \hsize{\dotfill}\medskip}

\newbox\kk

\begin{document}

\setbox\kk\vbox{\hsize=4.4cm\footnotesize\tabcolsep=2pt
\begin{center}
\begin{tabular}{|c|c|c|c|c|}
\hline
1 & 2 & 3 & $\ldots$ & 100 \\
\hline
101 & 102 & 103 & $\ldots$ & 200 \\
\hline
$\vdots$ & $\vdots$ & $\vdots$ & $\ddots$ & $\vdots$ \\
\hline
9901 & 9902 & 9903 & $\ldots$ & 10\,000 \\
\hline
\end{tabular}
\end{center}}

{\small
\centerline{\sc САНКТ-ПЕТЕРБУРГСКАЯ}
\centerline{\sc ОЛИМПИАДА ШКОЛЬНИКОВ 2024 года ПО МАТЕМАТИКЕ}
\centerline{\sc II тур. \ 9 класс.}
\medskip
\hrule
\bigskip

\q1. У Димы есть красный и синий фломастеры, одним из них он красит на числовой оси рациональные точки, а другим --- иррациональные. Дима покрасил 100 рациональных и 100 иррациональных точек, после чего стер подписи, позволявшие узнать, где находилось начало координат и какой был масштаб. У Сергея есть циркуль, которым он может замерить расстояние между любыми двумя покрашенными точками $A$ и~$B$, а~потом отметить на оси точку, находящуюся на замеренном расстоянии от любой покрашенной точки $C$ (слева или справа); при этом Дима тут же красит её соответствующим фломастером. Как Сергею узнать, каким цветом Дима красит рациональные точки, а каким --- иррациональные?

\q2. Силач Бамбула может одновременно нести несколько гирь, если их суммарный вес не превосходит 200~кг, и этих гирь не больше трёх. По дороге на работу он повредил палец и обнаружил, что теперь может нести не более двух гирь (и по-прежнему не более 200~кг). При каком наименьшем $k$ верно утверждение: \emph{любой набор из 100 гирь, которые Бамбула раньше мог перенести за 50 заходов, с больным пальцем он сможет перенести не более чем за $k$ заходов?}

\q3. Треугольник $ABC$ вписан в окружность. Из точек $B$ и $C$ одновременно выползают два муравья. Он ползут по дуге $BC$ навстречу друг другу так, что произведение расстояний от них до точки $A$ остается неизменным. Докажите, что во время их движения (до момента встречи) прямая, проходящая через муравьёв, касается некоторой фиксированной окружности.

\q4. Тренер выстроил в ряд 200 волейболистов и раздал им $m$
мячей (каждый волейболист мог получить сколько угодно
мячей). Время от времени один из волейболистов кидает мяч другому (а~тот ловит).
Через некоторое время оказалось, что из любых двух
волейболистов левый кидал правому мяч ровно два раза, а
правый левому ровно один раз. При каком наименьшем $m$ это
возможно?

\dotline

\centerline{Олимпиада 2024 года. II тур. 9 класс. Выводная аудитория.}
\smallskip

\q5. $AH$~--- высота остроугольного  треугольника $ABC$, вписанного в окружность $s$. На отрезке $BH$ выбрали точки $D$ и $E$. На лучах $AD$ и~$AE$ выбрали точки $X$ и $Y$ соответственно так, что середины отрезков $DX$ и $EY$ лежат на окружности $s$. Оказалось, что точки $B$, $X$, $Y$ и $C$ лежат на одной окружности. Докажите, что $BD+BE=2CH$.

\q6. Натуральное число $n$ назовём  {\it бедным}, если
уравнение
$$
x_1x_2\ldots x_{101}=(n-x_1)(n-x_2)\ldots(n-x_{101})
$$ 
не имеет решений в натуральных числах $1<x_i<n$. Существует ли бедное число, имеющее более
100\,000 различных простых делителей?

\q7. В очень большом Городе строят метро: много
станций, некоторые пары которых соединены тоннелями, причем от любой станции можно добраться по тоннелям до любой
другой. Все тоннели метро требуется разбить на <<линии>>:
каждая линия состоит из нескольких последовательных
тоннелей, все станции в которых различны (в~частности, линия не должна быть кольцевой); допускаются и линии, состоящие из одного тоннеля. По закону требуется, чтобы от любой станции
до любой другой можно было доехать, сделав не более 100
пересадок с линии на линию. При каком наибольшем $N$ любое
связное метро с $N$ станциями можно разбить на линии, соблюдая закон?
}


\clearpage

\centerline{\sc САНКТ-ПЕТЕРБУРГСКАЯ}
\centerline{\sc ОЛИМПИАДА ШКОЛЬНИКОВ 2024 года ПО МАТЕМАТИКЕ}
\centerline{\sc II тур. \ 10 класс.}
\medskip
\hrule
\medskip

\q1. В клетках доски $2024\times 2024$ расставлены целые числа так, что в 
любом прямоугольнике $2 \times 2023$ (вертикальном или горизонтальном) с одной вырезанной угловой клеткой, не выходящем за пределы доски, сумма чисел делится на 13. Докажите, что сумма всех чисел на доске делится на~13.

\q2. Даны одинаковые на вид 32 настоящие и 32 фальшивые монеты. Все фальшивые монеты весят поровну и меньше настоящих, которые тоже все весят одинаково. Как за шесть взвешиваний на весах с~двумя чашами определить тип хотя бы семи монет?

\q3. На стороне $BC$ остроугольного треугольника $ABC$ отмечена точка $P$. Точка $E$ симметрична точке $B$ относительно прямой~$AP$. Отрезок $PE$ пересекает описанную окружность треугольника $ABP$ в точке $D$. Точка $M$~--- середина стороны $AC$. Докажите, что ${DE+AC>2BM}$.

\q4. Рассмотрим всевозможные квадратные трехчлены вида ${x^2+ax+b}$, где 
$a$ и $b$~--- натуральные числа, не превосходящие некоторого натурального числа $N$. Докажите, что количество пар 
таких трехчленов, имеющих общий корень, не превосходит $N^2$.

\dotline

\centerline{Олимпиада 2024 года. II тур. 10 класс. Выводная аудитория.}
\smallskip

\q5. На числовой оси отмечено 2\,000\,000 точек с целыми координатами. Рассматриваются отрезки длин  97, 100 и 103 с концами в~этих точках. Какое наибольшее количество таких отрезков может быть?

\q6. Дан вписанный шестиугольник $AB_1CA_1BC_1$. Окружность~$\omega$ вписана и в треугольник $ABC$, и в треугольник $A_1B_1C_1$ и касается отрезков $AB$ и $A_1B_1$ в точках $D$ и $D_1$ соответственно. Докажите, что  если $\angle ACD = \angle BCD_1$, то $\angle A_1C_1D_1 = \angle B_1C_1D$.

\q7. Ребра полного графа на $1000$ вершинах покрашены в три цвета. Докажите, что в этом графе имеется несамопересекающийся одноцветный цикл, длина которого нечётна и не меньше 41.


\clearpage

\centerline{\sc САНКТ-ПЕТЕРБУРГСКАЯ}
\centerline{\sc ОЛИМПИАДА ШКОЛЬНИКОВ 2024 года ПО МАТЕМАТИКЕ}
\centerline{\sc II тур. \ 11 класс.}
\medskip
\hrule
\bigskip

\setrisbox\copy\kk

\q1. Таблица $100\times 100$ заполнена числами от 1 до 10\,000 так, как показано на рисунке.
Можно ли переставить некоторые числа так, чтобы 
\ris[\rismove(0pt,-4mm)]
в~каждой клетке по-прежнему стояло одно число, и чтобы во всех прямоугольниках из трех клеток сумма чисел не изменилась?

\q2. Дана последовательность $a_n$: 
$$
1, \ 2, \ 2, \ 3, \ 3, \ 3, \ 4, \ 4, \ 4, \ 4, \ \dots 
$$
\endris

\kern-4pt\noindent
(одна единица, две двойки, три тройки и т.\,д.) и еще одна последовательность $b_n$ такая, что $a\strut_{b_n}=b_{a_n}$ для всех натуральных $n$. Известно, что $b_{k}\!=\!1$ при некотором $k\!>\!100$. Докажите, что $b_m=1$ при всех $m>k$.

\q3. В неравнобедренном треугольнике $ABC$ проведена биссектриса~$AK$. Диаметр $XY$ его описанной окружности перпендикулярен прямой~$AK$ (порядок точек на описанной окружности $B$--$X$--$A$--$Y$--$C$). Окружность, проходящая через точки $X$ и $Y$, пересекает отрезки $BK$ и~$CK$ в точках $T$ и $Z$ соответственно. Докажите, что если $KZ = KT$, то $XT \perp YZ$.

\q4. Дано 101-значное число $a$ и произвольное натуральное число $b$. Докажите, что найдется такое не более чем 102-значное натуральное число $c$, что любое число вида $\overline{caaa\ldots ab}$ --- составное.

\dotline

\centerline{Олимпиада 2024 года. II тур. 11 класс. Выводная аудитория.}
\smallskip

\q5. На плоскости отмечено 100 точек общего положения (т.\,е.\ никакие три не лежат на одной прямой). Докажите, что можно выбрать три отмеченные точки $A, B, C$ так, чтобы для любой точки $D$ из оставшихся 97 отмеченных точек, прямые $AD$ и $CD$ не содержали бы точек, лежащих внутри треугольника $ABC$.

\q6. Дан многочлен $P(x)$ с целыми коэффициентами. Для некоторого натурального $n$ числа $P(0), P(1), \ldots , P(2^n + 1)$ делятся на $2^{2^n}$. Докажите, что значения многочлена $P(x)$ во всех целых точках делятся на $2^{2^n}$.

\q7. Турист прибыл на остров, где живут 100 волшебников, каждый из которых может быть рыцарем или лжецом. Он знает, что на момент его приезда один из ста волшебников~--- рыцарь (но не знает, кто именно), а~остальные~--- лжецы. Турист может выбрать любых двух волшебников $A$ и $B$ и попросить $A$ заколдовать $B$ заклинанием \emph{Вжух!}, которое меняет сущность (превращает рыцаря в лжеца, а лжеца в рыцаря). Волшебники выполняют просьбы туриста, но если в тот момент волшебник $A$ --- рыцарь, то сущность $B$ действительно меняется, а если $A$ --- лжец, то не меняется. Турист хочет после нескольких последовательных просьб одновременно знать сущность хотя бы $k$ волшебников. При каком наибольшем $k$ он сможет добиться своей цели?


\end{document}
